% GENERATED FILE. Edit canonical lessons, then run npm run generate:books.
% canonical-source-hash: fnv1a64:6585f3aa520edfab
% canonical-lessons: HI-C259-becaini, HI-C259-bhir, HI-C259-ghatna, HI-C259-durghatna, HI-C259-khatra

\chapter{Restlessness, A crowd, An event, An accident, Danger}
\label{ch:hi-restlessness-a-crowd-an-event-an-accident-danger}
\hlchaptermodality{259}

\begin{chapteropening}
\textbf{By the end of this chapter:} \emph{I can say restlessness, a crowd, an event, an accident and danger.}

Complete the last lesson of chapter 259: I can say restlessness, a crowd, an event, an accident and danger.
\end{chapteropening}

\section[\texorpdfstring{\hi{बेचैनी} (becainī)}{बेचैनी (becainī)}]{\hi{बेचैनी} (becainī) — restlessness, unease}

\label{lesson:HI-C259-becaini}



\begin{quote}
Before the new one: say the Hindi for faith, then the Hindi for loneliness.
\end{quote}

\subsection*{You'll want to know: \hi{बेचैनी}}
\textbf{\hi{बेचैनी}} — \emph{becainī} — \textquotedblleft{}restlessness, unease\textquotedblright{}.

\begin{cousinweb}
\textbf{\hi{बेचैनी}} (\emph{becainī}) is \emph{becain}, \textquotedblleft{}restless\textquotedblright{}, with the ending \emph{-ī} that makes a noun. \emph{be-} is the Persian \textquotedblleft{}without\textquotedblright{}; \emph{cain}, \textquotedblleft{}ease\textquotedblright{}, is of uncertain origin, and CDIAL has no entry for it.
\end{cousinweb}

\begin{grammarlens}[title={What you've built}]
One more word for this chapter.
\end{grammarlens}

\subsection*{Your turn}
\begin{itemize}
\raggedright
  \item \emph{Say it:} \emph{becainī}
  \item \emph{Say it:} \emph{becainī}, once more
  \item \emph{Say it:} say \emph{becainī}
  \item \emph{Recall:} say the Hindi for faith, then the Hindi for loneliness, then say \emph{becainī} again
\end{itemize}

\begin{tcolorbox}[breakable,colback=teal!4,colframe=teal!35!black,title={Before you move on}]
What does \hi{बेचैनी} mean? (\textquotedblleft{}Restlessness, unease\textquotedblright{}.) Say it once more.
\end{tcolorbox}
\section[\texorpdfstring{\hi{भीड़} (bhīṛ)}{भीड़ (bhīṛ)}]{\hi{भीड़} (bhīṛ) — a crowd}

\label{lesson:HI-C259-bhir}



\begin{quote}
Before the new one: say the Hindi for loneliness, then the Hindi for restlessness.
\end{quote}

\subsection*{You'll want to know: \hi{भीड़}}
\textbf{\hi{भीड़}} — \emph{bhīṛ} — \textquotedblleft{}a crowd\textquotedblright{}.

\begin{cousinweb}
\textbf{\hi{भीड़}} (\emph{bhīṛ}) is a tadbhava, a word inherited and reshaped by sound change. It comes from \emph{*bhiṭ}, \textquotedblleft{}meet, throng\textquotedblright{}, a form CDIAL reconstructs from its descendants (CDIAL 9490).
\end{cousinweb}

\begin{grammarlens}[title={What you've built}]
One more word for this chapter.
\end{grammarlens}

\subsection*{Your turn}
\begin{itemize}
\raggedright
  \item \emph{Say it:} \emph{bhīṛ}
  \item \emph{Say it:} \emph{bhīṛ}, once more
  \item \emph{Say it:} say \emph{bhīṛ}
  \item \emph{Recall:} say the Hindi for loneliness, then the Hindi for restlessness, then say \emph{bhīṛ} again
\end{itemize}

\begin{tcolorbox}[breakable,colback=teal!4,colframe=teal!35!black,title={Before you move on}]
What does \hi{भीड़} mean? (\textquotedblleft{}A crowd\textquotedblright{}.) Say it once more.
\end{tcolorbox}
\section[\texorpdfstring{\hi{घटना} (ghaṭnā)}{घटना (ghaṭnā)}]{\hi{घटना} (ghaṭnā) — an event}

\label{lesson:HI-C259-ghatna}



\begin{quote}
Before the new one: say the Hindi for restlessness, then the Hindi for a crowd.
\end{quote}

\subsection*{You'll want to know: \hi{घटना}}
\textbf{\hi{घटना}} — \emph{ghaṭnā} — \textquotedblleft{}an event\textquotedblright{}.

\begin{cousinweb}
\textbf{\hi{घटना}} (\emph{ghaṭnā}) is Sanskrit \emph{ghaṭanā}, \textquotedblleft{}a joining, a coming about\textquotedblright{}, taken over whole as a learned loan, a tatsama. It is a different word from \emph{ghaṭnā}, \textquotedblleft{}to decrease\textquotedblright{} (CDIAL 4415).
\end{cousinweb}

\begin{grammarlens}[title={What you've built}]
One more word for this chapter.
\end{grammarlens}

\subsection*{Your turn}
\begin{itemize}
\raggedright
  \item \emph{Say it:} \emph{ghaṭnā}
  \item \emph{Say it:} \emph{ghaṭnā}, once more
  \item \emph{Say it:} say \emph{ghaṭnā}
  \item \emph{Recall:} say the Hindi for restlessness, then the Hindi for a crowd, then say \emph{ghaṭnā} again
\end{itemize}

\begin{tcolorbox}[breakable,colback=teal!4,colframe=teal!35!black,title={Before you move on}]
What does \hi{घटना} mean? (\textquotedblleft{}An event\textquotedblright{}.) Say it once more.
\end{tcolorbox}
\section[\texorpdfstring{\hi{दुर्घटना} (durghaṭnā)}{दुर्घटना (durghaṭnā)}]{\hi{दुर्घटना} (durghaṭnā) — an accident}

\label{lesson:HI-C259-durghatna}



\begin{quote}
Before the new one: say the Hindi for a crowd, then the Hindi for an event.
\end{quote}

\subsection*{You'll want to know: \hi{दुर्घटना}}
\textbf{\hi{दुर्घटना}} — \emph{durghaṭnā} — \textquotedblleft{}an accident\textquotedblright{}.

\begin{cousinweb}
\textbf{\hi{दुर्घटना}} (\emph{durghaṭnā}) is Sanskrit \emph{durghaṭanā}, \textquotedblleft{}a bad occurrence\textquotedblright{}, taken over whole as a learned loan, a tatsama. It is \emph{ghaṭanā}, \textquotedblleft{}event\textquotedblright{}, with \emph{dur-}, \textquotedblleft{}bad\textquotedblright{}.
\end{cousinweb}

\begin{grammarlens}[title={What you've built}]
One more word for this chapter.
\end{grammarlens}

\subsection*{Your turn}
\begin{itemize}
\raggedright
  \item \emph{Say it:} \emph{durghaṭnā}
  \item \emph{Say it:} \emph{durghaṭnā}, once more
  \item \emph{Say it:} say \emph{durghaṭnā}
  \item \emph{Recall:} say the Hindi for a crowd, then the Hindi for an event, then say \emph{durghaṭnā} again
\end{itemize}

\begin{tcolorbox}[breakable,colback=teal!4,colframe=teal!35!black,title={Before you move on}]
What does \hi{दुर्घटना} mean? (\textquotedblleft{}An accident\textquotedblright{}.) Say it once more.
\end{tcolorbox}
\section[\texorpdfstring{\hi{ख़तरा} (khatrā)}{ख़तरा (khatrā)}]{\hi{ख़तरा} (khatrā) — danger}

\label{lesson:HI-C259-khatra}



\begin{quote}
Before the new one: say the Hindi for an event, then the Hindi for an accident.
\end{quote}

\subsection*{You'll want to know: \hi{ख़तरा}}
\textbf{\hi{ख़तरा}} — \emph{khatrā} — \textquotedblleft{}danger\textquotedblright{}.

\begin{cousinweb}
\textbf{\hi{ख़तरा}} (\emph{khatrā}) is a loan from Arabic \emph{khaṭar}, \textquotedblleft{}danger, risk\textquotedblright{}, which reached Hindi through Persian. Its Arabic root is \emph{kh-ṭ-r} (Lane).
\end{cousinweb}

\begin{grammarlens}[title={What you've built}]
One more word for this chapter.
\end{grammarlens}

\subsection*{Your turn}
\begin{itemize}
\raggedright
  \item \emph{Say it:} \emph{khatrā}
  \item \emph{Say it:} \emph{khatrā}, once more
  \item \emph{Say it:} say \emph{khatrā}
  \item \emph{Recall:} say the Hindi for an event, then the Hindi for an accident, then say \emph{khatrā} again
\end{itemize}

\begin{tcolorbox}[breakable,colback=teal!4,colframe=teal!35!black,title={Before you move on}]
What does \hi{ख़तरा} mean? (\textquotedblleft{}Danger\textquotedblright{}.) Say it once more.
\end{tcolorbox}
